\section{From a coefficient relation to a planar surface}
\label{sec:topology}

We now prove Theorem~\ref{thm:main}. The construction has two steps. A
finite product of conjugates of the relator gives a punctured disk with
paired bands. We then remove its connected band surfaces from the inside
out, applying Theorem~\ref{thm:planar} at each step.
The relation-to-surface interpretation is also described in
\cite[preprint Section 2, especially Example 2.5]{Chen2026}. We retain the finite
domain and the original coefficient map explicitly in the construction.

\begin{proof}[Proof of Theorem~\ref{thm:main}]
Cyclically conjugating the relator does not change its normal closure.
Since its total $t$-exponent is $\pm1$, it has a cyclic expression beginning
with a $t$-letter. Splitting powers of $t$ into individual letters gives
\eqref{eq:word}; identity coefficients are allowed. Suppose that
$g\in A$ belongs to the normal closure of $w$ in $A*\langle t\rangle$.
Then, for some finite $r$,
\begin{equation}\label{eq:normal-closure-product}
 g=\prod_{j=1}^{r} b_jw^{\varepsilon_j}b_j^{-1},
 \qquad b_j\in A*\langle t\rangle,\quad
 \varepsilon_j\in\{1,-1\}.
\end{equation}
If $r=0$, then $g=1$. Henceforth assume $r>0$.

\paragraph{Realizing the relation.}
Choose a based CW presentation space $Y$ with $\pi_1(Y)=A$. Such a space
exists for every group: take one loop $x_a$ for each $a\in A$ and attach
two-cells imposing $x_1=1$ and $x_ax_b=x_{ab}$. These relations identify
its fundamental group with $A$; see also
\cite[Corollary 1.28]{Hatcher2002}. Thus
$\pi_1(Y\vee\T)=A*\langle t\rangle$, with the added circle representing
$t$. All subsequent triangulations concern a finite planar domain.

Let $\Delta$ be a polygonal disk, and remove the interiors of $r$ disjoint
polygonal disks to obtain a planar surface $P$. There is a continuous map
\[
 f:P\longrightarrow Y\vee\T
\]
whose outer boundary is a based loop in $Y$ representing $g$, and whose
$j$th inner boundary reads $w^{\varepsilon_j}$ when oriented as the
boundary of its missing disk. To construct it, join marked points on the
inner circles to a common outer basepoint by cutting arcs that are otherwise
disjoint. Order the arcs so that the outer loop represents the ordered
product of the conjugated missing-disk loops, and map the $j$th arc to a
path representing $b_j$. Use chosen based coefficient loops in $Y$ and
one monotone traversal of $\T$, affine in angle, for each $t^{\pm1}$;
on an inverse relator, reverse the chosen paths in reverse order.
Cutting along the arcs gives a
disk whose boundary path is nullhomotopic by
\eqref{eq:normal-closure-product}. Extend over that disk and reglue the
matching copies of the arcs. This uses the equality in
$\pi_1(Y\vee\T)$, before imposing the relator $w$.

Write $u:P\to\T$ for the projection of $f$ that collapses $Y$, and
$v:P\to Y$ for the projection that collapses $\T$. The boundary map $u$
is constant on the outer circle and on every coefficient interval, and
traverses $\T$ once on each $t$-letter interval.

\paragraph{Extracting bands while retaining the original map.}
Identify $\T$ with $\mathbb R/\Z$, with wedge point $0$. Approximate
$u$, relative to the boundary, by a map $u'$ that is piecewise affine in
local angle coordinates, with uniform circle distance less than $1/12$.
Here is a direct construction. Subdivide the boundary so that its angle
maps are affine on each edge, and choose a sufficiently fine finite
triangulation of $P$. By uniform continuity, the image under $u$ of each
triangle lies in an arbitrarily short arc. Lift to angles on each triangle
and interpolate the vertex values affinely. The lifts on shared edges
differ by an integer, so the resulting circle maps agree there. On boundary
edges the interpolation equals $u$. Finer triangulations make the uniform
error arbitrarily small. Only the compact surface $P$ and its circle map
have been triangulated.

Choose an angle $z\in(1/3,2/3)$ outside the images of all vertices. The
level $u'^{-1}(z)$ is a finite disjoint union of polygonal circles and
properly embedded polygonal arcs: within each triangle it is a line
segment or empty, and it misses every vertex. Because $u'=u$ on the
boundary, each $t$-letter interval contains exactly one arc endpoint;
there are no other endpoints. Ignore the closed level circles.

The circle orientation coorients the level arcs. Their endpoints have
opposite crossing signs along the induced boundary orientation of $P$.
The missing-disk orientation reverses this boundary orientation at both
ends, so each arc pairs opposite $t$-letter signs in the word-disk
readings. This remains true for an arc with both ends on the same disk.

At every point of the level, the original value $u$ has distance greater
than $1/4$ from the wedge point. Consequently the original map $f$ takes
an open neighborhood of the level into $\T\setminus\{0\}$, and $v$ is
identically the basepoint there. Take pairwise disjoint thin rectangular
neighborhoods of the level arcs inside this neighborhood. They give bands
meeting the missing-disk boundaries in small intervals inside the
corresponding $t$-letter intervals. Thus $v$ is exactly constant throughout
every band.

Each small attachment interval is now a slot, with the sign of its
original $t$-letter. Between consecutive slots, the projected boundary
path consists of the original coefficient loop together with constant
pieces from the adjacent $t$-intervals. Its label in $A$ is therefore
unchanged. In particular positive disks have the corners $a_i$, and
negative disks have the reversed corners $a_i^{-1}$. This observation
allows the geometric attachment slots to be smaller than the original
letter intervals.

\paragraph{Removing the planar components.}
Put the missing disks back into $\Delta$ as word disks and adjoin all
the bands. Their union is a finite compact planar surface $S$, possibly
disconnected. Each component contains a word disk, every slot is used
once, and the paired signs are opposite. Its orientation is induced from
the plane, so it satisfies the disk-and-band hypotheses of
Theorem~\ref{thm:planar}. Along its boundary, $v$ reads exactly the
corner words from that theorem: band sides and remaining pieces of
$t$-intervals are constant in $Y$.

The nonexterior boundary words of an arbitrary component need not yet be
trivial, because other components may lie inside them. We therefore choose
an innermost component as follows. Every connected planar component $C$
has a unique exterior boundary circle; let $D(C)$ be the closed Jordan
disk it bounds. Its other boundary circles bound the inner complementary
disks of $C$. Among the finitely many $D(C)$ choose one minimal under
inclusion. No inner complementary disk of this $C$ contains another
component: such a component would have its exterior disk strictly inside
$D(C)$, contradicting minimality.

Each inner complementary disk of $C$ therefore contains no word disks.
It already carries the map $v$ to $Y$, so its boundary word is trivial in
$A$. Theorem~\ref{thm:planar} makes the exterior boundary word trivial
as well. Hence its actual boundary loop under $v$ is nullhomotopic in
$Y$. Fill all of $D(C)$ by a map to $Y$ agreeing with that boundary loop,
and discard the word disks and bands of $C$.

We continue to write $v$ for the resulting coefficient map.
This replacement changes nothing outside $D(C)$ and affects no other
component. It leaves a continuous $Y$-map outside the remaining word-disk
interiors, with the same corner readings and constant remaining bands.
That is precisely the invariant needed to repeat the argument; no further
circle approximation is required. Each step removes a component. After
finitely many steps, the original outer loop extends over $\Delta$ in
$Y$. Thus $g=1$ in $A$, proving the claimed injectivity.
\end{proof}
